\chapter{Problema de pesquisa e objetivos}\label{problema-de-pesquisa-e-objetivos}

\section{Problema de pesquisa}\label{problema-de-pesquisa}

Dados AMI possibilitam observar alterações que não aparecem no faturamento agregado, mas a mesma riqueza de dados também cria novas fontes de ambiguidade. Uma anomalia pode existir no consumo, na qualidade de comunicação, na relação entre fases, no comportamento temporal ou no próprio cadastro do equipamento. Além disso, a população de consumidores é heterogênea e está distribuída em alimentadores com características distintas.

O problema científico não é, portanto, apenas separar valores normais de valores raros. É necessário \textbf{localizar a discrepância na UC e no período, caracterizar sua natureza e distinguir evidência de anomalia de evidência de causa}. O problema operacional acrescenta uma segunda restrição: mesmo quando muitas UCs apresentam desvios, somente uma parcela pode ser inspecionada.

A pergunta principal do trabalho é:

\begin{quote}
\textbf{Como detectar, localizar e caracterizar padrões anômalos em dados AMI de baixa frequência, combinando evidências temporais, elétricas, comportamentais, de comunicação e contextuais, com ênfase em padrões compatíveis com perdas não técnicas?}
\end{quote}

Dessa pergunta derivam duas questões auxiliares:

\begin{enumerate}
\def\labelenumi{\arabic{enumi}.}
\item
  \textbf{Como decompor as diferentes fontes de evidência sem confundir falhas de telemetria, eventos de rede, mudanças legítimas de comportamento e potenciais irregularidades?}
\item
  \textbf{Como transformar as evidências produzidas pelo modelo em uma priorização rastreável e compatível com a capacidade hipotética de inspeção?}
\end{enumerate}

\section{Objetivo geral}\label{objetivo-geral}

Desenvolver, implementar e avaliar retrospectivamente uma metodologia e uma plataforma de software para \textbf{detecção, localização e caracterização de padrões anômalos em redes de distribuição a partir de dados AMI}, combinando engenharia de dados temporal, técnicas estatísticas, aprendizado de máquina e contexto elétrico/cadastral. A finalidade é produzir evidências interpretáveis e listas analíticas de candidatos relevantes para investigação de perdas não técnicas, sem pressupor confirmação de fraude.

\section{Objetivos específicos}\label{objetivos-especuxedficos}

\begin{enumerate}
\def\labelenumi{\arabic{enumi}.}
\tightlist
\item
  Integrar dados AMI, cadastrais e topológicos em arquitetura com identidade e vigência temporal explícitas.
\item
  Reconstruir vínculos históricos UC--NIO, tratando trocas de equipamento, cadências e registradores sem induzir descontinuidades espúrias.
\item
  Implementar extração incremental, processamento em lotes, persistência, auditoria de integridade e construção rastreável dos datasets.
\item
  Extrair indicadores de comunicação, consumo, forma diária, mudança de regime, grandezas elétricas e eventos disponíveis.
\item
  Implementar e examinar detectores não supervisionados, regras determinísticas e especialistas, registrando suas saídas e critérios de combinação.
\item
  Analisar retrospectivamente, na população estudada, efeitos de elegibilidade, mudança de ranking e composição de listas sob diferentes parâmetros e variantes metodológicas.
\item
  Demonstrar a exploração dos dados, os modelos e as hipóteses diagnósticas por meio de uma TUI, do protótipo analítico e de estudos de caso auditáveis.
\item
  Formular cenários de aproveitamento das listas analíticas sob restrições de capacidade e um protocolo de avaliação posterior com registros independentes de inspeção.
\end{enumerate}

Os sete primeiros objetivos são examinados por implementações, experimentos e resultados computacionais. O oitavo compreende \textbf{projeção e desenho de avaliação}, não mensuração de rendimento de vistorias: o estudo não executou inspeções nem recebeu registros de OS ou TOI.

\section{Rastreabilidade dos objetivos às evidências}\label{rastreabilidade-dos-objetivos-uxe0s-eviduxeancias}

As evidências do trabalho são separadas pela natureza do resultado, evitando equiparar um escore computacional a um resultado de fiscalização.

\begin{longtable}[]{@{}
  >{\raggedright\arraybackslash}p{(\columnwidth - 4\tabcolsep) * \real{0.3333}}
  >{\raggedright\arraybackslash}p{(\columnwidth - 4\tabcolsep) * \real{0.3333}}
  >{\raggedright\arraybackslash}p{(\columnwidth - 4\tabcolsep) * \real{0.3333}}@{}}
\caption{Correspondência entre objetivos e evidências}\label{tab:integracao1}\\
\toprule\noalign{}
\begin{minipage}[b]{\linewidth}\raggedright
Objetivos
\end{minipage} & \begin{minipage}[b]{\linewidth}\raggedright
Evidências incluídas
\end{minipage} & \begin{minipage}[b]{\linewidth}\raggedright
Alcance da conclusão
\end{minipage} \\
\midrule\noalign{}
\endfirsthead
\toprule\noalign{}
\begin{minipage}[b]{\linewidth}\raggedright
Objetivos
\end{minipage} & \begin{minipage}[b]{\linewidth}\raggedright
Evidências incluídas
\end{minipage} & \begin{minipage}[b]{\linewidth}\raggedright
Alcance da conclusão
\end{minipage} \\
\midrule\noalign{}
\endhead
\bottomrule\noalign{}
\endlastfoot
1--3 & Arquitetura UC--NIO, campanha de extração, particionamento, auditoria e datasets & Implementação e execução do pipeline \\
4--5 & Features, detectores, normalização, gating, fusão, casos e métricas de ranking & Comportamento analítico e reprodutibilidade retrospectiva \\
6 & Série \texttt{DAILY\_RULES}, inferência \texttt{OOT\_MULTIVIEW}, ablação e sensibilidade & Volume de listas, interseções e estabilidade; não acerto de fraude \\
7 & TUI, protótipo de Cockpit e figuras de casos AMI & Visualização e rastreabilidade das evidências disponíveis \\
8 & Cenário de capacidade e formulações econômicas; protocolo de campo & Possibilidade de aplicação futura, sem vistorias realizadas \\
\end{longtable}
{\small\noindent Fonte: elaboração própria.\par}


As contagens só são comparadas dentro de coortes, cortes e políticas identificados.

\section{Hipótese de pesquisa}\label{hipuxf3tese-de-pesquisa}

A hipótese investigada é que a decomposição da análise em qualidade, evidências elétricas, comportamento e contexto torna os resultados mais rastreáveis e altera de forma identificável a priorização computacional em relação à abordagem global. Os ensaios retrospectivos permitem examinar essas propriedades e a sensibilidade das filas a decisões metodológicas. Uma hipótese distinta --- de que a política elevaria a taxa de acerto na identificação de perdas não técnicas --- exige resultados independentes de vistoria e não é avaliada neste TCC.
